% Partie : méthodes
% Chapitre : démonstrations
% Section : premier exemple

\documentclass[../../../include/open-logic-section]{subfiles}

\begin{document}
\olfileid[fr]{mth}{prf}{ex1}

\olsection{Un exemple}

Notre premier exemple est un fait simple concernant
l'union et l'intersection d'ensembles. Il illustrera
le dépliage des définitions, la démonstration de
conjonctions et d'énoncés universels, ainsi que
la démonstration par cas.

\begin{prop}
Pour tous ensembles $A$, $B$ et $C$,
$A \cup (B \cap C) = (A \cup B) \cap (A \cup C)$.
\end{prop}

Démontrons-le !{}

\begin{proof}
Nous voulons montrer que, pour tous ensembles
$A$, $B$ et $C$,
$A \cup (B \cap C) = (A \cup B) \cap (A \cup C)$.
\begin{quote}
Commençons par déplier, dans la proposition,
la définition de «~$=$~». Rappelons que,
pour démontrer que deux ensembles sont
identiques, il faut montrer qu'ils ont
les mêmes !!{element}s. Autrement dit,
tous les !!{element}s de $A \cup (B \cap C)$
doivent aussi être !!{element}s de
$(A \cup B) \cap (A \cup C)$, et
réciproquement. «~Réciproquement~» signifie
que tout !!{element} de
$(A \cup B) \cap (A \cup C)$ doit aussi
être !!a{element} de $A \cup (B \cap C)$.
En dépliant la définition, nous découvrons
donc une conjonction à démontrer.
Notons-la :
\end{quote}
Par définition,
$A \cup (B \cap C) = (A \cup B) \cap (A \cup C)$
si et seulement si tout !!{element} de
$A \cup (B \cap C)$ est !!a{element} de
$(A \cup B) \cap (A \cup C)$, et si tout
!!{element} de $(A \cup B) \cap (A \cup C)$
est !!a{element} de $A \cup (B \cap C)$.
\begin{quote}
Puisqu'il s'agit d'une conjonction,
il faut démontrer séparément chacun
de ses membres. Commençons par le
premier : tout !!{element} de
$A \cup (B \cap C)$ est
!!a{element} de
$(A \cup B) \cap (A \cup C)$.

Cet énoncé est universel. Considérons
donc un !!{element} arbitraire de
$A \cup (B \cap C)$ et montrons qu'il
est aussi !!a{element} de
$(A \cup B) \cap (A \cup C)$.
Choisissons une variable pour nommer
cet !!{element} arbitraire,
par exemple~$z$. Poursuivons :
\end{quote}
Montrons d'abord que tout !!{element} de
$A \cup (B \cap C)$ est
!!a{element} de
$(A \cup B) \cap (A \cup C)$.
Soit $z \in A \cup (B \cap C)$.
Nous devons montrer que
$z \in (A \cup B) \cap (A \cup C)$.
\begin{quote}
Il est temps de déplier les définitions
de $\cup$ et de~$\cap$. Par exemple,
la définition de l'union dit que
$A \cup B =
  \Setabs{u}{u \in A \text{ ou } u \in B}$.
Pour l'appliquer à «~$A \cup (B \cap C)$~»,
il faut remplacer le «~$B$~» de la
définition par «~$B \cap C$~». Ainsi,
$A \cup (B \cap C) =
  \Setabs{u}{u \in A \text{ ou } u \in B \cap C}$.
Notre hypothèse $z \in A \cup (B \cap C)$
signifie donc que
$z \in
  \Setabs{u}{u \in A \text{ ou } u \in B \cap C}$.
Or $z \in \Setabs{u}{\dots u\dots}$
si et seulement si $\dots z\dots$.
Dans le cas présent, cela signifie
$z \in A$ ou $z \in B \cap C$.
\end{quote}
Par définition de $\cup$, on a $z \in A$
ou $z \in B \cap C$.
\begin{quote}
Nous avons une disjonction : il sera
utile de distinguer deux cas.
Dans chacun d'eux, nous montrerons
la conclusion visée, à savoir
«~$z \in (A \cup B) \cap (A \cup C)$~».
\end{quote}
Cas 1 : supposons $z \in A$.
\begin{quote}
Cette hypothèse ne nous donne guère
d'autres informations. Regardons donc
la conclusion à établir. Pour montrer
que $z \in (A \cup B) \cap (A \cup C)$,
la définition de $\cap$ nous demande
de montrer que $z$ appartient à la
fois à $A \cup B$ et à $A \cup C$.
Or $z \in A \cup B$ si et seulement
si $z \in A$ ou $z \in B$, et nous
avons déjà $z \in A$ par l'hypothèse
du premier cas. Le même raisonnement,
avec $C$ à la place de $B$, donne
$z \in A \cup C$. Nous avons raisonné
à rebours à partir de la conclusion ;
écrivons maintenant l'argument dans
le sens requis par la démonstration.
\end{quote}
Puisque $z \in A$, nous avons
$z \in A$ ou $z \in B$ ; par
définition de~$\cup$, $z \in A \cup B$.
De même, $z \in A \cup C$. Par
définition de~$\cap$, il en résulte
$z \in (A \cup B) \cap (A \cup C)$.
\begin{quote}
Le premier cas est terminé. Passons
au second, dans lequel $z \in B \cap C$.
\end{quote}
Cas 2 : supposons $z \in B \cap C$.
\begin{quote}
Il s'agit à nouveau de l'intersection
de deux ensembles. Appliquons la
définition de~$\cap$ :
\end{quote}
Puisque $z \in B \cap C$, $z$ est
!!a{element} de $B$ et de $C$,
par définition de~$\cap$.
\begin{quote}
Revenons à la conclusion : il faut
montrer que $z$ appartient à la
fois à $A \cup B$ et à $A \cup C$.
La solution est, là encore, immédiate.
\end{quote}
Puisque $z \in B$, nous avons
$z \in A \cup B$. Puisque
$z \in C$, nous avons aussi
$z \in A \cup C$. Donc
$z \in (A \cup B) \cap (A \cup C)$.
\begin{quote}
Nous venons d'appliquer de nouveau
les définitions de $\cup$ et de $\cap$.
Comme nous les avons déjà rappelées
et avons déjà montré qu'un objet
appartenant à l'un de deux ensembles
appartient à leur union, nous pouvons
être moins explicites cette fois.

Le second cas est achevé : nous pouvons
énoncer la conclusion de la première
moitié de la démonstration.
\end{quote}
Ainsi, si $z \in A \cup (B \cap C)$,
alors $z \in (A \cup B) \cap (A \cup C)$.
\begin{quote}
Il reste à montrer l'autre inclusion :
tout !!{element} de
$(A \cup B) \cap (A \cup C)$ est
!!a{element} de $A \cup (B \cap C)$.
Comme précédemment, nous traitons cet
énoncé universel en prenant un
!!{element} arbitraire du premier
ensemble et en montrant qu'il appartient
au second. Énonçons cette démarche.
\end{quote}
Supposons maintenant
$z \in (A \cup B) \cap (A \cup C)$.
Nous voulons montrer
$z \in A \cup (B \cap C)$.
\begin{quote}
Notre nouvelle hypothèse est
$z \in (A \cup B) \cap (A \cup C)$.
Nous réutilisons ici la lettre~$z$,
ce qui peut prêter à confusion.
Mais les hypothèses posées dans
la première moitié ne sont plus
en vigueur : nous pouvons maintenant
en poser de nouvelles à propos de~$z$.
Si vous préférez, remplacez $z$ par
une autre variable dans la suite.

Par définition de $\cap$, $z$ appartient
à $A \cup B$ ainsi qu'à $A \cup C$.
En dépliant $\cup$, nous obtenons
$z \in A$ ou $z \in B$, et aussi
$z \in A$ ou $z \in C$. On pourrait
penser à distinguer des cas, mais
le «~et~» entre les deux disjonctions
demande de l'attention. On pourrait
croire à tort que les seules
possibilités sont simplement :
$z$ appartient à $A$, à $B$ ou
à $C$. Examinons plutôt chaque
disjonction tour à tour.
\end{quote}
Par définition de $\cap$,
$z \in A \cup B$ et
$z \in A \cup C$.
Par définition de $\cup$,
$z \in A$ ou $z \in B$.
Distinguons deux cas.
\begin{quote}
Nous nous concentrons sur la
première disjonction et n'avons
pas encore déplié la seconde,
issue de $A \cup C$ : nous
n'en avons pas encore besoin.
Le premier cas est $z \in A$.
Tout !!a{element} d'un ensemble
est aussi !!a{element} de l'union
de cet ensemble avec un autre.
Ce cas est donc facile :
\end{quote}
Cas 1 : supposons $z \in A$.
Alors $z \in A \cup (B \cap C)$.
\begin{quote}
Dans le second cas, $z \in B$.
Nous allons y déplier la
seconde union et distinguer
de nouveau des cas :
\end{quote}
Cas 2 : supposons $z \in B$.
Puisque $z \in A \cup C$,
$z \in A$ ou $z \in C$.
Distinguons encore deux cas.

Cas 2a : $z \in A$. Alors, de nouveau,
$z \in A \cup (B \cap C)$.
\begin{quote}
Ce cas peut sembler étrange :
nous n'avons pas eu besoin
de l'hypothèse $z \in B$.
Cela ne pose aucun problème.
\end{quote}
Cas 2b : $z \in C$. Alors
$z \in B$ et $z \in C$, donc
$z \in B \cap C$ et, par suite,
$z \in A \cup (B \cap C)$.
\begin{quote}
Les deux distinctions de cas
sont achevées, ainsi que la
seconde moitié de la démonstration.
\end{quote}
Ainsi, si
$z \in (A \cup B) \cap (A \cup C)$,
alors $z \in A \cup (B \cap C)$.
\end{proof}

\end{document}
